\documentclass{situationist-preprint}

\title{Errata to \enquote{On Free Will}, Volume II}
\msid{ACT-01-e}
\msfamily{ACT-01}
\msdiscipline{Action, Agency and Free Will}
\mstype{Erratum}
\mskeywords{corrigendum; consequence argument; Rule Beta; agglomeration; reasons-responsiveness}

\begin{document}
\maketitle

\begin{abstract}
We correct six errors in Volume~II of ACT-01 (\emph{Determinism and the Consequence Argument}, Chapters 11--20): three misnumbered references, a misprint in the statement of Rule Beta, a lemma stated without a necessary side-condition, and a citation error. For each we give the printed text, the correction, and the results that depend on it. No theorem of the volume changes its statement except Lemma~15.2, whose corrected form still supports every result derived from it. The conclusion of the work, that free will is compatible both with determinism and with indeterminism, is unaffected.
\end{abstract}

Volume~II defends the first half of the independence assumption of Volume~I (ACT-01-a, Assumption~3.1) against the consequence argument, and argues that the argument fails at its transfer rule. Pagination is continuous across the three volumes; Volume~II occupies pp.~323--641, and page numbers below are those of the printed work.

\section{Misnumbered references}

\paragraph{1. Chapter~11, p.~331, line~4.} The opening summary of Volume~I states that the compatibility thesis \enquote{follows from Assumption~3.2 of Volume~I}. Assumption~3.2 does not exist; the independence assumption is Assumption~3.1, and Theorem~3.2 is the compatibility thesis that follows from it. The sentence should read \enquote{follows from Assumption~3.1 of Volume~I}. The same reference is printed correctly on p.~334 and in the proof of Theorem~13.4.

\paragraph{2. Chapter~16, proof of Theorem~16.3 (p.~497).} Theorem~16.3 states that the original Rule Beta is invalid. Its proof derives Agglomeration from Rule Alpha and Rule Beta \enquote{by Proposition~14.3}. The derivation is Proposition~14.5; Proposition~14.3 is the validity of Rule Alpha on every reading of \enquote{no choice} considered in Chapter~14, which the derivation uses but which does not itself yield Agglomeration. The reference should read \enquote{by Proposition~14.5}. The proof is otherwise correct.

\paragraph{3. Chapter~19, p.~602, note~11.} The note states that reasons-responsiveness is defined in \enquote{Definition~2.3 of Volume~I}. Definition~2.3 of Volume~I defines ownership of a mechanism; reasons-responsiveness is Definition~2.2. The note should read \enquote{Definition~2.2 of Volume~I}. The argument of the surrounding section uses reasons-responsiveness, as the text makes clear, and is unaffected.

\section{A misprint in the statement of Rule Beta}

\paragraph{4. Chapter~12, p.~341, display~(12.2).} Rule Beta is printed as: from $Np$ and $N(q \to p)$ infer $Nq$. The second premise should read $N(p \to q)$. The rule as printed is not the one discussed anywhere in the volume; every subsequent statement of it, including those on pp.~342, 389 and 497, and the derivation of the consequence argument in display~(12.4), uses the correct form.

\section{A missing side-condition}

\paragraph{5. Chapter~15, Lemma~15.2 (p.~452).} Chapter~15 presents the coin counterexample to Agglomeration due to \citet{McKayJohnson1996}. Lemma~15.2 is printed as follows.

\begin{quote}\small
\textbf{Lemma 15.2.} If an agent could have tossed a fair coin at $t$ and did not, she has no choice about whether the coin lands heads at $t$, and no choice about whether it lands tails at $t$.
\end{quote}

As stated, the lemma is false of an agent who could toss the coin in a way that determines its outcome, for instance by a practiced technique that reliably produces heads. Such an agent has a choice about whether the coin lands heads, since her tossing, done in the right way, would make it land so. The lemma needs the side-condition that the outcome of a toss, were she to make one, would not vary with the agent's reasons. The corrected statement is:

\begin{quote}\small
\textbf{Lemma 15.2 (corrected).} If an agent could have tossed a fair coin at $t$ and did not, and if, had she tossed it, which way it landed would not have varied with her reasons, then she has no choice about whether the coin lands heads at $t$, and no choice about whether it lands tails at $t$.
\end{quote}

The condition is satisfied in the case described in Chapter~15, where the coin is tossed by a mechanism that the agent can start but not influence, and in every application of the lemma in Chapters~15 and~16. In particular, the counterexample to Agglomeration (Proposition~15.3) and the invalidity of the original Rule Beta (Theorem~16.3) stand as printed. The side-condition is also the one that the account of choice in Chapter~14 makes natural: on that account, having a choice about $p$ requires that $p$'s truth would vary with one's reasons through one's own mechanism of action, which is the reasons-responsiveness of Volume~I, Definition~2.2.

\section{A citation error}

\paragraph{6. Chapter~17, p.~538, and bibliography.} \citet{Huemer2000} is listed as appearing in \emph{Mind}. It appeared in \emph{The Philosophical Review}, 109(4). The text's report of the paper, as a discussion of whether a version of the consequence argument can avoid the McKay--Johnson counterexample, is correct. The neighboring citation of \citet{Lewis1981} on the same page, for the distinction between a strong and a weak sense in which one might be able to break a law, is correct as printed.

\section{Effect on the work}

None of the corrections changes the statement of a theorem or proposition. Items~1, 2, 3 and~6 affect references only. Item~4 is a misprint in a display that is not used. Item~5 adds to a lemma a condition that every application of it already satisfies. The volume's thesis, that the consequence argument fails at its transfer rule, and with it the work's conclusion that free will is compatible both with determinism and with indeterminism (ACT-01-a, Theorem~3.2), are unaffected.

\begin{thebibliography}{9}
\bibitem[Huemer(2000)]{Huemer2000} Huemer, M. (2000). Van Inwagen's consequence argument. \emph{The Philosophical Review}, 109(4), 525--544.
\bibitem[Lewis(1981)]{Lewis1981} Lewis, D. (1981). Are we free to break the laws? \emph{Theoria}, 47(3), 113--121.
\bibitem[McKay and Johnson(1996)]{McKayJohnson1996} McKay, T.~J. and Johnson, D. (1996). A reconsideration of an argument against compatibilism. \emph{Philosophical Topics}, 24(2), 113--122.
\end{thebibliography}

\end{document}
